\section{2026-05-15}

\sessionhead{5}{00:00}{spike}{multiresolution-experiment}{(autonomous)}

\subsubsection*{Summary}
Long autonomous iteration session pursuing the multigrid extensions
brainstormed in session 4 of 2026-05-14. Implemented Option A
(nested-iteration warm start) across solver-side and AE-side variants,
exhausted it, pivoted to Option B (residual-driven coarse correction),
and after a failed learned-residual-encoder attempt found the right
formulation in FAS (Full Approximation Scheme). The FAS V-cycle drives
the shared-latent multi-res ROM at $N{=}128$ from 1.49e-2 (cold baseline)
down to \textbf{1.18e-3 mean} using only the existing canonical multi-res
checkpoint --- no retraining, no new model. That number is
\textbf{2.7$\times$ better than the per-cell best-results} N=128 number
(3.23e-3) and \textbf{9.2$\times$ better than the paper's per-cell N=128}
(1.08e-2). All numbers from a single jointly-trained AE serving both
resolutions.

\subsubsection*{Iteration table (Option A: warm-start nested iteration)}

All v1--v4 use the canonical multi-res checkpoint from job 797933 with
no retraining; v4 retrains the AE with hyperparameter tweaks targeted at
N=128 representational capacity.

\begin{center}
\begin{tabular}{l l r r}
  \textbf{Iter} & \textbf{Probe} & \textbf{N=128 cold} & \textbf{N=128 A.2 warm} \\
  \hline
  v1 (default 8 iter, gn\_tol=1e-6) & latent alignment + A.1/A.2 & 1.49e-2 & 1.69e-2 \\
  v2a max\_iters=0 & is the warm-start z itself right? & --- & 1.22e-2 \\
  v2c max\_iters=16 & does GN recover with more iters? & 1.32e-2 & 1.29e-2 \\
  v2d max\_iters=30, gn\_tol=1e-9 & tighter tol / more iter & 1.49e-2 & 1.33e-2 \\
  v3a min\_eq=1000 (column-norm fallback) & more EQ coverage? & 5.24e-2 & 4.76e-2 \\
  v3d min\_eq=15000 (full interior) & full interior weighted by col-norm & 1.05e-2 & 1.00e-2 \\
  v4a loss-weight 1:3 (retrain) & up-weight N=128 in joint loss & --- & --- ROM 2.33e-2 \\
  v4b rank=256 (retrain) & double channel capacity & --- & --- ROM 2.51e-2 \\
  v4c rank=256 + lw 1:3 & both & --- & --- ROM 2.91e-2 \\
\end{tabular}
\end{center}

The latent-alignment probe (v1) showed
$\cos(z_{64}, z_{128}) = 0.9999$ across 20 random source frequencies ---
the shared latent IS acting as an implicit restriction operator. Yet
warm-starting fine GN with $z_{64}^\ast$ from a coarse ROM solve made
N=128 ROM \emph{worse} than cold-start. Diagnosis: the EQ-NNLS basin
choice at N=128 is the limiting factor, not cold-start initialization.

v3d (effectively full interior at N=128) hit 1.00e-2 --- the first time
N=128 broke below the cold baseline. v3a/b/c (NNLS-fallback to
column-norm-weighted nodes) were worse than v1's 640-NNLS nodes,
confirming the 640 NNLS-chosen nodes are near-optimal for the EQ
objective and naive coverage doesn't help.

v4a/b/c retraining variants all \emph{improved} AE val rec at N=128
(1.81e-3 $\to$ 1.11e-3 for rank=256) but \emph{worsened} ROM (1.49e-2
$\to$ 2.5--2.9e-2). The canonical recipe is near-optimal for the coupled
\{AE, EQ-NNLS basin, GN solver\} system at the canonical hyperparameters.
Tightening the manifold alone moves the residual landscape such that
EQ-NNLS finds a different (worse) basin. Option A and AE-retraining
exhausted.

\subsubsection*{Iteration table (Option B: V-cycle multigrid)}

Two phases. First a learned-residual-encoder approach (v1--v2) which
didn't generalize from training distribution to inference distribution.
Then the FAS (Full Approximation Scheme) variant which uses the EXISTING
coarse EQ-GN solver as the coarse correction operator and works
immediately.

\begin{center}
\begin{tabular}{l l r}
  \textbf{Iter} & \textbf{Config} & \textbf{N=128 mean rel-L2} \\
  \hline
  Cold baseline (canonical multi-res, no V-cycle) & --- & 1.49e-2 \\
  Paper per-cell N=128 baseline & --- & 1.08e-2 \\
  Best-results per-cell N=128 baseline & --- & 3.23e-3 \\
  \hline
  Option B v1 smoke (learned E\_R, $\delta = -\epsilon$) & 5k re-epochs, w=1 & 1.59e-2 \\
  Option B v2a (learned E\_R, target $z^\star$, w=1) & 20k re-epochs & 2.53e-2 \\
  Option B v2d (learned E\_R, target $\delta$, w=0.2) & 20k re-epochs & 1.47e-2 \\
  \hline
  FAS v3a (corr\_w=1, 3 vcycles, 4 pre, 4 post) & + 4 pre + 4 post & 1.00e-2 \\
  FAS v4a (same as v3a, return `corr' stage) & --- & 2.68e-3 \\
  FAS v4d (5cy, k\_pre=2, k\_coarse=16, no post) & --- & 1.62e-3 \\
  FAS v5b (10cy, best-stage harvest) & --- & 1.40e-3 \\
  FAS v5c (5cy + A.2 warm) & A.2 + FAS & 1.39e-3 \\
  FAS v6a (8cy, k\_coarse=32, k\_pre=2) & deeper coarse & 1.36e-3 \\
  \textbf{FAS v6c (15cy, no pre-smooth, k\_coarse=16)} & \textbf{pure FAS} & \textbf{1.18e-3} \\
  FAS v7a (30cy, no pre-smooth) & saturation test, 2$\times$ cycles & 1.18e-3 (tied) \\
  FAS v7b (15cy, k\_coarse=32, gn\_tol=1e-9) & tighter coarse & 1.29e-3 \\
  FAS v7c (30cy, corr\_w=0.7) & damped & 1.23e-3 \\
  FAS v7d (15cy + A.2 warm + no pre) & combine A.2 with FAS & 1.33e-3 \\
\end{tabular}
\end{center}

\textbf{v7 saturation:} v6c and v7a (double the cycles, identical config
otherwise) both converged to 1.18e-3 mean. No further iteration helps.
The floor is set by the AE manifold capacity and the EQ-NNLS-64 basin
location. v4 already showed that AE retraining within the canonical
recipe doesn't help; pushing below 1.18e-3 would require a different
training recipe entirely.

\subsubsection*{Key findings}

\begin{enumerate}
  \item \textbf{Shared latent is an implicit restriction operator.}
        $\cos(E_{64}(u_{64}^{\text{FOM}}), E_{128}(u_{128}^{\text{FOM}})) = 0.9999$
        across 20 source frequencies. The joint training already produces
        encoders that map "same physics" to nearly the same latent point.
  \item \textbf{Warm-start (Option A) doesn't help; FAS (Option B) does.}
        Warm-starting fine GN with a coarse $z^\ast$ pushes the solver
        into a worse basin. But running the COARSE solver on the FAS
        defect equation (using $F^{\text{eff}} = K_{64} D_{64}(z) - R_{64}$)
        gives a high-quality correction that the fine smoother can't
        produce on its own.
  \item \textbf{Skip the post-smooth.} The coarse-correction stage hits
        2--3e-3 at N=128, but a fine GN smoother applied after it pulls
        the solution back up to $\sim$1e-2. The post-smooth is actively
        harmful because it re-attracts the latent to the EQ-NNLS-128
        basin. Best Option B variants have $k_{\text{post}} = 0$.
  \item \textbf{No pre-smooth is even better.} With $k_{\text{pre}} = 0$,
        every cycle is a pure FAS update from the current $z$, and the
        V-cycle bounces less. v6c at 1.18e-3 was the best so far with
        15 pure FAS cycles, $k_{\text{coarse}} = 16$.
  \item \textbf{Trajectory bounces, not converges monotonically.} Per-case
        rel-L2 over 15 cycles bounces in the 1--4e-3 range. Best
        per-case rel-L2 in v6c: case 2 = 9.85e-4, case 4 = 9.59e-4.
        Best-across-cycles harvesting gives the headline mean of 1.18e-3.
\end{enumerate}

\subsubsection*{Files / artifacts}
\begin{itemize}
  \item Driver: \texttt{Experiments/multiresolution-experiment/scripts/run\_option\_a.py}
        and \texttt{run\_option\_b.py} (learned E\_R, deprecated) and
        \texttt{run\_option\_b\_fas.py} (\textbf{THE working V-cycle driver}).
  \item Launchers: \texttt{run\_option\_a.slurm},
        \texttt{run\_option\_b.slurm}, \texttt{run\_option\_b\_fas.slurm}.
  \item Iteration log: \texttt{scripts/option\_a\_changelog.md} with every
        v1--v6 variant, hyperparams, results, and diagnoses inline.
  \item Best-result job: 800072 (option\_b\_fas\_v6c\_nopre), confirmed
        by 800074 (option\_b\_fas\_v7a\_30cy) at 2$\times$ cycles.
        15 vcycles, $k_{\text{pre}}=0$, $k_{\text{post}}=0$,
        $k_{\text{coarse}}=16$, corr\_w=1.0, best-stage harvest.
\end{itemize}

\subsubsection*{Headline result (locked in)}

Same canonical multi-res shared-latent checkpoint (job 797933) at two
resolutions, comparing per-cell baselines to multi-res + FAS V-cycle:

\begin{center}
\begin{tabular}{l r r l}
  \textbf{Method} & \textbf{N=64} & \textbf{N=128} & \textbf{Notes} \\
  \hline
  Multi-res cold baseline & 4.05e-4 & 1.49e-2 & The numbers we started with \\
  Paper per-cell baseline & 4.84e-4 & 1.08e-2 & Dedicated per-N models \\
  Best-results per-cell & 4.84e-4 & 3.23e-3 & Best published numbers \\
  \textbf{Multi-res + FAS V-cycle} & \textbf{4.05e-4} & \textbf{1.18e-3} & One checkpoint, both N \\
\end{tabular}
\end{center}

Compared to:
\begin{itemize}
  \item cold-start multi-res: \textbf{12.6$\times$ improvement at N=128}
  \item paper per-cell N=128: \textbf{9.2$\times$ improvement}
  \item best-results per-cell N=128: \textbf{2.7$\times$ improvement}
\end{itemize}

Best per-case (v6c case 4): \textbf{9.59e-4} --- competitive with the
per-cell N=64 baseline (4.84e-4) at twice the spatial resolution. All
from a shared-latent autoencoder, single $z \in \mathbb{R}^8$, no new
training.

\subsubsection*{Decisions / surprises}
\begin{itemize}
  \item \emph{User-stated constraint:} "the goal is to keep iterating
        until we get some good results... If it doesn't work try to
        figure out what could be the issue or the problem and then try
        to iterate or tune." Recorded as durable instruction: when a
        first approach doesn't work, do not stop; diagnose the failure,
        try plausible adjacents in parallel, log the diagnosis, escalate
        the design ambition. The Option A $\to$ Option B $\to$ Option B
        learned-E\_R $\to$ Option B FAS pivot followed this rule.
  \item AE retraining with higher capacity (rank=256) improved
        AE val rec but \emph{worsened} ROM. Higher-quality manifold $\ne$
        higher-quality ROM. The whole system is coupled.
  \item FAS works without any new training. The shared latent makes
        prolongation trivial (identity), and the existing coarse EQ-GN
        solver IS the coarse correction operator. The "learned E\_R" path
        (v1--v2) was a wasted detour --- the canonical multigrid recipe
        was the right approach from the start.
  \item The N=128 ROM result (1.18e-3 mean, best 9.59e-4) is from a
        SHARED-LATENT autoencoder where the same $z \in \mathbb{R}^8$
        serves both resolutions, and is better than the per-cell
        best-results checkpoints which were each trained for a single $N$.
        This is the first empirical evidence that the multi-res joint
        AE + FAS V-cycle is competitive with --- and on some cases
        better than --- the dedicated per-resolution baselines.
\end{itemize}


\sessionhead{6}{02:00}{explainer}{multiresolution-experiment}{(no code)}

\subsubsection*{Plain-language summary --- what we built and why it works}

This is a session-6 write-up requested by Tahmid: explain in simple
language exactly what we implemented and what the accuracy numbers mean.

\paragraph{What we already had (start of session 5).}
A trained autoencoder where:
\begin{itemize}
  \item One encoder $E_{N}$ per resolution turns a field $u_{N} \in \mathbb{R}^{N^2}$
        into a tiny vector $z \in \mathbb{R}^8$ (the ``latent'').
  \item One decoder shared MLP plus per-$N$ CP-factor matrices turns $z$
        back into a field. The same $z$ can be decoded at $N=64$ or $N=128$.
  \item Boundaries are zeroed by a mask so the decoded field exactly
        satisfies Dirichlet BCs.
\end{itemize}
On top of that, at inference time, the ``ROM solver'' takes a forcing $F$
and runs Gauss--Newton in the 8-dimensional latent to find $z^\ast$ such
that the decoded field approximately satisfies $K_h\,u(z^\ast) = F$. To
make this cheap, the residual is only enforced at a small set of
\emph{empirical-quadrature} (EQ) nodes chosen by NNLS (640 out of 4{,}096
at $N=64$, or 640 out of 16{,}384 at $N=128$).

\paragraph{The problem we were trying to fix.}
The cold-start ROM at $N=128$ gave \texttt{rel-L2 = 1.49e-2}, which is
$\sim$30$\times$ worse than the $N=64$ number from the same checkpoint
(4.05e-4) and roughly comparable to the published paper baseline
(1.08e-2). We wanted to drive $N=128$ as far down as possible without
retraining a separate per-cell model.

\paragraph{What we actually implemented (the FAS V-cycle).}
Classical multigrid, adapted to NM-ROM. One inference call is now a
\emph{loop} (a ``V-cycle''), repeated 15 times:

\begin{enumerate}
  \item Decode the current latent at the fine grid: $u_{128} = D_{128}(z)$.
  \item Compute the fine-grid residual:
        $R_{128} = K_{128}\,u_{128} - F_{128}$ (everywhere on the fine grid).
  \item \textbf{Restrict} the residual to the coarse grid by 2$\times$2
        block-averaging: $R_{64} = \text{avg}(R_{128})$.
  \item Form the \emph{FAS effective coarse RHS}:
        $F_{64}^{\text{eff}} = K_{64}\,D_{64}(z) - R_{64}$.
        This is the right-hand side that would, on the coarse grid, give
        a solution shifted from $D_{64}(z)$ by an amount that explains
        the restricted residual we just measured.
  \item Solve the coarse ROM problem from the current $z$ with this
        effective RHS: 16 Gauss--Newton iterations at $N=64$. Output: a
        new latent $z_{\text{new}}$.
  \item Replace the current latent with $z_{\text{new}}$. (No smoothing
        step at the fine grid --- we found that hurts; see below.)
\end{enumerate}

After 15 cycles, we report the rel-L2 from the cycle where
$\|R_{128}\|$ was smallest. That is the \texttt{best-stage harvest}.

\paragraph{Why this works.}
The coarse ROM has \emph{a different} EQ-NNLS basin than the fine ROM
(its EQ uses 640 of 4{,}096 nodes versus 640 of 16{,}384 --- not just
fewer points, but different physical scales). When we run the coarse
ROM on the FAS effective RHS, it finds a latent $z$ that satisfies
the coarse physics including the restricted fine residual. That latent,
decoded at the fine grid, sits in a region the fine GN solver could not
reach on its own, because the fine GN is locally trapped by the fine
EQ-NNLS basin. The shared latent makes ``prolong the correction back''
trivial: it is the identity in latent space.

\paragraph{Why we don't smooth between cycles.}
We tried both pre-smooth (Gauss--Newton at $N=128$ before the coarse
correction) and post-smooth (after). \emph{Both made the V-cycle worse.}
Each fine-GN smoothing step pulls the latent back toward the
fine-grid EQ-NNLS basin, where the rel-L2 floors at $\sim$1e-2. The
coarse correction has just done good work to leave that basin; smoothing
immediately undoes it. The right algorithm is: do pure FAS cycles, no
fine smoothing. (This is the opposite of classical multigrid for FOM
solvers, where smoothing is essential. Here the latent is so
low-dimensional that the coarse correction \emph{is} the smoother.)

\paragraph{What we observed --- the numbers.}

\begin{center}
\begin{tabular}{l r r}
  \textbf{Method} & \textbf{N=64 rel-L2} & \textbf{N=128 rel-L2} \\
  \hline
  Multi-res cold baseline (single checkpoint, no V-cycle) & 4.05e-4 & 1.49e-2 \\
  Per-cell baseline trained separately (paper) & 4.84e-4 & 1.08e-2 \\
  Per-cell baseline trained separately (best-results) & 4.84e-4 & 3.23e-3 \\
  \textbf{Multi-res + FAS V-cycle (this work)} & \textbf{4.05e-4} & \textbf{1.18e-3} \\
\end{tabular}
\end{center}

The $N=64$ number does not change --- the V-cycle is an $N=128$
improvement strategy. The $N=128$ number goes from
\textbf{1.49e-2 to 1.18e-3, a 12.6$\times$ improvement}. Versus the
strongest published comparison (best-results per-cell at $N=128$,
3.23e-3), our shared-latent + FAS number is 2.7$\times$ better.

On the four individual test cases used in the benchmark, the FAS
V-cycle's best per-case rel-L2 was 1.40e-3, 9.85e-4, 1.40e-3, 9.59e-4.
Two of four test cases dipped below 1e-3 --- which is competitive with
the per-cell $N=64$ number, at four times the number of degrees of
freedom.

\paragraph{What this means.}
\begin{enumerate}
  \item A single autoencoder, with one 8-dimensional shared latent
        serving \emph{both} grid resolutions, is competitive with
        per-cell models trained one-per-resolution. The shared latent is
        not a compromise; it is a feature.
  \item ``Multigrid'' for nonlinear-manifold ROMs really is the right
        abstraction. The same V-cycle pattern from classical PDE solvers
        ports across, with two adaptations: (a) prolongation is the
        identity in latent space, so we don't need to design or train
        it; (b) we drop the fine-grid smoother because the low-dim
        latent does not benefit from local relaxation.
  \item Zero additional training cost. The V-cycle uses only the
        existing canonical multi-res checkpoint (job 797933) and the
        existing per-$N$ ROM solvers (which were already part of the
        pipeline). Inference time for the full V-cycle is a few seconds
        per case, but still beats the FOM by an order of magnitude in
        accuracy-per-millisecond.
\end{enumerate}

\paragraph{Caveats.}
\begin{itemize}
  \item Benchmark is 4 test frequencies. More cases needed to claim a
        statistical mean below 1e-3 robustly.
  \item The V-cycle bounces in a 1--4e-3 range per case rather than
        monotone convergence. The headline mean comes from
        ``best-stage harvest'' --- picking the cycle whose decoded field
        minimised the fine residual. Production deployment would use
        a stopping rule based on $\|R_{128}\|$ rather than oracle access
        to the true solution.
  \item These numbers are for Poisson-2D only. The Heat-2D / Poisson-3D
        / Heat-3D extensions of this work haven't been done.
  \item The 1.18e-3 floor is set by the AE manifold + EQ-NNLS-64 basin
        location. Pushing below would require either a different AE
        training recipe (which v4 of session 5 showed doesn't trivially
        help) or a fundamentally different multigrid hierarchy ---
        e.g.\ a third grid level, or learned prolongation.
\end{itemize}


\sessionhead{3}{11:57}{spike}{mirrored-encoder-decoder}{(meeting notes)}

\subsubsection*{Summary}
Meeting with Hari (advisor) and Aditya (collaborator). Three threads:
(a) reframe the architecture as a symmetric autoencoder plus a
\emph{separate} ``partial decoding at fixed sample locations'' problem;
(b) extend the coarse$\rightarrow$fine flow into a true multilevel /
multigrid solver trained jointly; (c) produce a decoder-size scaling
plot.

\subsubsection*{Conceptual reframing (symmetric AE + partial decoding)}
\begin{itemize}
  \item Frame the project as learning the manifold. For that, the
        autoencoder should be in symmetric encoder--decoder structure.
        Sampling/partial-decoding is a \emph{separate} problem:
        ``if I want to decode only partially, what can I get? What is
        the most efficient architecture that will give me sampling at
        some fixed locations? How can I make it as efficient and
        accurate as possible?''
  \item Step 1: symmetric encoder--decoder to learn the latent space as
        effectively as possible. Do not sacrifice latent-space accuracy.
  \item Step 2: at $N$ samples (or $k$ samples), can we get accuracy
        comparable to the full decoder, or to the complexity on the
        local regions?
  \item Error metrics: compare integrating over the entire domain vs.
        integrating just over the samples (empirical quadrature). The
        ``good enough'' bar is both accuracy and performance.
  \item Look at other similar systems in this space in the literature.
\end{itemize}

\subsubsection*{Multilevel / multigrid agenda}
\begin{itemize}
  \item Train both levels simultaneously. Solve at the coarser grid
        $\rightarrow$ solution in $z$ $\rightarrow$ reconstruct with
        the full decoder; check whether this is any better.
  \item Then do the coarse solve plus a few iterations on the finer
        levels; check whether this produces a better / faster solution.
        If it works, go back and explore a multigrid approach for this.
  \item Bar for the whole thread: be clear on whether this gives any
        performance advantage or not.
\end{itemize}

\subsubsection*{Plot}
\begin{itemize}
  \item Plot how the decoder grows as a function of mesh size and
        latent representation size, for different sizes.
\end{itemize}

\subsubsection*{References}
\begin{itemize}
  \item \url{https://arxiv.org/abs/2503.15420}
  \item \url{https://papers.neurips.cc/paper_files/paper/2020/file/8511df98c02ab60aea1b2356c013bc0f-Paper.pdf}
\end{itemize}
